\pagestyle{empty}
\tight
\begin{tblr}{width=\textwidth,colspec={|Q[c,m,1.9cm]|X[l,m]|},hlines,vlines,hspan=minimal,row{1}={bg=headblue},rowsep=1pt,colsep=3pt}
\SetCell{c,m}\hfil\logo\hfil & \textbf{POLITEKNIK NEGERI MALANG}\par\textbf{JURUSAN TEKNOLOGI INFORMASI}\par\textbf{PROGRAM STUDI: D4 TEKNIK INFORMATIKA} \\
\SetCell[c=2]{c} \textbf{RENCANA PEMBELAJARAN SEMESTER (RPS)} & \\
\end{tblr}
\begin{longtblr}{width=\textwidth,colspec={|Q[l,m,1.9cm]|X[0.85,l,m]|X[1.45,l,m]|X[0.75,l,m]|X[1.15,l,m]|},hlines,vlines,hspan=minimal,rowsep=1pt,colsep=2pt,row{1,3}={bg=headgray,font=\bfseries}}
MATA KULIAH & KODE & BOBOT (SKS)/jam & SEMESTER & TGL. PENYUSUNAN \\
Aljabar Linier & RTI252002 & 2 SKS/4 jam & 2 & 29 Januari 2026 \\
OTORISASI & Dosen Pengembang RPS & Koordinator RMK & \SetCell[c=2]{l} Ka PRODI &  \\
 & Muhammad Afif Hendrawan, S.Kom., M.T. &  & \SetCell[c=2]{l}{\vspace{8mm}\par Dr. Ely Setyo Astuti, S.T., M.T.} &  \\
\SetCell[r=8]{l} Capaian Pembelajaran (CP) & \SetCell[c=4]{l,bg=headgray,font=\bfseries} Capaian Pembelajaran Lulusan Program Studi (CPL-Prodi) &  &  &  \\
 & CPL09 & \SetCell[c=3]{l} Memiliki pemahaman yang memadai terkait cara kerja sistem komputer dan berbagai teknik/pendekatan berbasis komputasi untuk memecahkan masalah stakeholder. &  &  \\
 & \SetCell[c=4]{l,bg=headgray,font=\bfseries} Capaian Pembelajaran mata kuliah (CPL-MK) &  &  &  \\
 & CPMK0902 & \SetCell[c=3]{l} Mahasiswa mampu menguasai prinsip matematis dan struktur diskrit serta menerapkan pendekatan komputasional sebagai dasar analisis terhadap permasalahan. &  &  \\
 & \SetCell[c=4]{l,bg=headgray,font=\bfseries} Kemampuan akhir tiap tahapan belajar (Sub-CPMK) &  &  &  \\
 & SCPMK0902-01101 & \SetCell[c=3]{l} Mahasiswa mampu menyelesaikan sistem persamaan linier menggunakan metode eliminasi dan representasi matriks, serta memvalidasi hasil secara komputasional. &  &  \\
 & SCPMK0902-01102 & \SetCell[c=3]{l} Mahasiswa mampu menerapkan operasi matriks (invers, transpose, determinan) untuk menyelesaikan permasalahan. &  &  \\
 & SCPMK0902-01103 & \SetCell[c=3]{l} Mahasiswa mampu menerapkan konsep vektor dan ruang vektor untuk interpretasi geometri dan komputasi dasar pada konteks informatika. &  &  \\
\SetCell[r=2]{l} Penilaian & \SetCell[c=4]{l} *Nilai CPMK didapat dari agregasi nilai SUB-CPMK yang dibebankan &  &  &  \\
 & \SetCell[c=4]{l}{\tiny\begin{tblr}{width=\linewidth,colspec={|Q[c,m,0.1\linewidth]|Q[c,m,0.16\linewidth]|X[l,m]|X[l,m]|X[l,m]|X[l,m]|X[l,m]|Q[c,m,0.08\linewidth]|},hlines,vlines,hspan=minimal,rowsep=1pt,colsep=0.7pt,row{1,2}={bg=headgray,font=\bfseries}}
Kode CPMK & Kode SCPMK & \SetCell[c=5]{c} Bobot per Bentuk Penilaian & & & & & Total Bobot \\
 & & Tugas & Kuis 1 & UTS & Kuis 2 & UAS & \\
\SetCell[r=1]{c} \SetCell[r=3]{c} CPMK0902 & SCPMK0902-01101 & 5\%: worksheet SPL mg 1--3 dan studi kasus & 10\%: pemodelan dan penyelesaian SPL & 12.5\%: SPL dan eliminasi Gauss/Gauss-Jordan & & 10\%: SPL dan validasi hasil & 37.5\% \\
\SetCell[r=2]{c}  & SCPMK0902-01102 & 9\%: worksheet matriks mg 5--7, 9 dan studi kasus & & 12.5\%: operasi matriks, invers, determinan & 5\%: aturan Cramer dan operasi matriks & 12.5\%: matriks, determinan, aturan Cramer & 39\% \\
 & SCPMK0902-01103 & 6\%: worksheet vektor mg 10--14 dan studi kasus & & & 5\%: vektor dot/cross product & 12.5\%: vektor, eigen, proyeksi ortogonal & 23.5\% \\
\SetCell[c=7]{r}\textbf{Total Per Penilaian} & & & & & & & \textbf{100\%} \\
\end{tblr}} &  &  &  \\
Diskripsi Singkat Mata Kuliah & \SetCell[c=4]{l} Mata kuliah Aljabar Linier membekali mahasiswa Teknik Informatika dengan konsep dan keterampilan dasar untuk memodelkan serta menyelesaikan masalah komputasi menggunakan pendekatan matematis. Materi mencakup sistem persamaan linier dan metode eliminasi (Gauss dan Gauss-Jordan), operasi matriks, invers dan transpos, determinan (ekspansi kofaktor/Laplace dan Sarrus), aturan Cramer untuk kasus kecil, operasi vektor di 2D/3D (dot product dan cross product), proyeksi ortogonal, serta eigenvalue dan eigenvector. Pembelajaran dirancang berorientasi Outcome-Based Education (OBE) dan student-centered learning melalui latihan terstruktur, pembahasan masalah, serta kasus-kasus sederhana yang relevan dengan bidang informatika sebagai fondasi untuk mata kuliah lanjutan seperti pembelajaran mesin, visi komputer, dan analisis data. &  &  &  \\
Bahan Kajian / Materi Pembelajaran & \SetCell[c=4]{l}\begin{enumerate}\item Permodelan Sistem Persamaan Linier (SPL)\item Notasi matriks SPL dan dasar operasi baris elementer\item Metode Gauss dan Gauss-Jordan\item Operasi matriks\item Inversi matriks\item Determinan\item Aturan Cramer\item Vektor\item Eigenvector dan eigenvalue\item Proyeksi ortogonal\end{enumerate} &  &  &  \\
\SetCell[r=2]{l} Pustaka & \SetCell[c=4]{l} \textbf{Utama:}\begin{enumerate}\item Lay, David C., et al. Linear Algebra and Its Applications. 6th ed., Pearson, 2021.\end{enumerate} &  &  &  \\
 & \SetCell[c=4]{l} \textbf{Pendukung:}\begin{enumerate}\item Hefferon, Jim. Linear Algebra. Orthogonal Publishing L3c, 2017.\end{enumerate} &  &  &  \\
Media Pembelajaran & \SetCell[c=2]{l} \textbf{Software:} LMS Polinema dan perangkat lunak komputasi numerik (opsional). &  & \SetCell[c=2]{l} \textbf{Hardware:} Komputer, laptop, proyektor. &  \\
Nama Dosen Pengampu & \SetCell[c=4]{l}\begin{enumerate}\item Muhammad Afif Hendrawan, S.Kom., M.T.\item Prof. Dr. Eng. Cahya Rahmad, S.T., M.T.\item Erfan Rohadi, S.T., M.Eng., Ph.D.\item Mungki Astiningrum, S.T., M.Kom.\item Triana Fatmawati, S.T., M.T.\item Adevian Fairuz Pratama, S.ST., M.Eng.\end{enumerate} &  &  &  \\
Matakuliah Syarat & \SetCell[c=4]{l} - &  &  &  \\
\end{longtblr}
\begin{longtblr}{width=\textwidth,colspec={|Q[c,m,0.06\textwidth]|X[1.35,l,m]|X[1.08,l,m]|X[1.16,l,m]|X[0.7,l,m]|X[1.24,l,m]|X[1.06,l,m]|X[0.98,l,m]|Q[c,m,0.05\textwidth]|},hlines,vlines,hspan=minimal,rowhead=2,row{1,2}={bg=headgreen,font=\bfseries},rowsep=1pt,colsep=1.5pt}
Mgg.\par Ke & Kemampuan Akhir Yang Direncanakan (Sub-CP-MK) & Bahan kajian (Materi Pembelajaran) & Bentuk dan Metode Pembelajaran & Estimasi Waktu & Pengalaman Belajar Mahasiswa & Kriteria \& Bentuk Penilaian & Indikator Penilaian & Bobot Penilaian (\%) \\
(1) & (2) & (3) & (4) & (5) & (6) & (7) & (8) & (9) \\
1 & SCPMK0902-01101 & Pengantar dan permodelan sistem persamaan linier (SPL). & Bentuk: kuliah luring. Metode: Self-Directed Learning. Penugasan: Latihan 1 permodelan dan penyelesaian SPL sederhana; Worksheet 1 penyelesaian SPL dengan metode eliminasi dan substitusi. & 2 x 2 x 50' tatap muka dan latihan terstruktur; 2 x 2 x 50' tugas mandiri (worksheet). & Mahasiswa membuat dan menyelesaikan sistem persamaan linier dari suatu masalah dengan metode eliminasi dan substitusi. & Kriteria: ketepatan dan penguasaan. Bentuk penilaian: ketepatan jawaban tugas. & Mampu membuat dan menyelesaikan sistem persamaan linier dari suatu masalah dengan metode eliminasi dan substitusi. & 1 \\
2 & SCPMK0902-01101 & Notasi matriks pada SPL; dasar operasi baris elementer / row echelon form (REF). & Bentuk: kuliah luring. Metode: Self-Directed Learning. Penugasan: Latihan 1 membuat notasi matriks dari SPL; Worksheet 1 menyelesaikan SPL dengan operasi baris elementer. & 2 x 2 x 50' tatap muka dan latihan terstruktur; 2 x 2 x 50' tugas mandiri (worksheet). & Mahasiswa memodelkan permasalahan SPL dalam bentuk matriks, memahami proses REF, dan menyelesaikan SPL dengan REF. & Kriteria: ketepatan dan penguasaan. Bentuk penilaian: ketepatan jawaban tugas. & Mampu memodelkan permasalahan SPL dalam bentuk matriks; memahami proses REF; menyelesaikan permasalahan SPL dengan REF. & 1 \\
3 & SCPMK0902-01101 & Penyelesaian SPL dengan metode Gauss (REF) dan metode Gauss-Jordan / reduced row echelon form (RREF). & Bentuk: kuliah luring. Metode: Self-Directed Learning. Penugasan: Latihan 1 dan Worksheet 1 menyelesaikan masalah SPL dengan metode Gauss dan Gauss-Jordan. & 2 x 2 x 50' tatap muka dan latihan terstruktur; 2 x 2 x 50' tugas mandiri (worksheet). & Mahasiswa mendapatkan nilai dari himpunan penyelesaian sistem persamaan linier dengan menggunakan metode Gauss dan Gauss-Jordan. & Kriteria: ketepatan dan penguasaan. Bentuk penilaian: ketepatan jawaban tugas. & Mampu menyelesaikan permasalahan SPL dengan metode Gauss (REF) dan Gauss-Jordan (RREF). & 2 \\
4 & SCPMK0902-01101 & Kuis 1: materi pekan 1--3. & Ujian tertulis. & Sesuai jadwal asesmen. & Mahasiswa mengerjakan soal kuis dengan materi pekan 1--3 secara mandiri. & Ujian. & Ketepatan jawaban dengan kunci jawaban. & 10 \\
5 & SCPMK0902-01102 & Operasi matriks: penjumlahan, pengurangan, perkalian skalar, dan perkalian antar matriks. & Bentuk: kuliah luring. Metode: Self-Directed Learning. Penugasan: Latihan 1 dan Worksheet 2 menyelesaikan soal operasi matriks. & 2 x 2 x 50' tatap muka dan latihan terstruktur; 2 x 2 x 50' tugas mandiri (worksheet). & Mahasiswa melakukan operasi penjumlahan, pengurangan, perkalian skalar, dan perkalian antar matriks. & Kriteria: ketepatan dan penguasaan. Bentuk penilaian: ketepatan jawaban tugas. & Mampu menjumlahkan dan mengurangi antar matriks; mengalikan matriks dengan bilangan skalar; mengalikan antar matriks. & 2 \\
6 & SCPMK0902-01102 & Inversi matriks dan partisi matriks. & Bentuk: kuliah luring. Metode: Self-Directed Learning. Penugasan: Latihan 1 dan Worksheet 2 menyelesaikan soal inversi matriks dan partisi matriks. & 2 x 2 x 50' tatap muka dan latihan terstruktur; 2 x 2 x 50' tugas mandiri (worksheet). & Mahasiswa melakukan proses inversi matriks ukuran $n \times n$ serta membuat partisi matriks dan melakukan operasi dan inversi pada partisi matriks. & Kriteria: ketepatan dan penguasaan. Bentuk penilaian: ketepatan jawaban tugas. & Mampu melakukan proses inversi matriks; membuat partisi matriks; melakukan operasi dan inversi pada partisi matriks. & 2 \\
7 & SCPMK0902-01102 & Determinan dan ekspansi kofaktor. & Bentuk: kuliah luring. Metode: Self-Directed Learning. Penugasan: Latihan 1 menyelesaikan soal determinan; Worksheet 3 determinan matriks $n \times n$. & 2 x 2 x 50' tatap muka dan latihan terstruktur; 2 x 2 x 50' tugas mandiri (worksheet). & Mahasiswa mencari nilai determinan matriks ukuran $n \times n$ dengan menggunakan ekspansi kofaktor. & Kriteria: ketepatan dan penguasaan. Bentuk penilaian: ketepatan jawaban tugas. & Mampu menghitung nilai determinan dengan menggunakan ekspansi kofaktor. & 2 \\
8 & SCPMK0902-01101, SCPMK0902-01102 & Ujian Tengah Semester: materi pekan 1--7. & Ujian tertulis. & Sesuai jadwal asesmen. & Mahasiswa mengerjakan soal ujian dengan materi pekan 1--7 secara mandiri. & Ujian. & Ketepatan jawaban dengan kunci jawaban. & 25 \\
9 & SCPMK0902-01102 & Aturan Cramer. & Bentuk: kuliah luring. Metode: Self-Directed Learning. Penugasan: Latihan 1 dan Worksheet 3 menyelesaikan kasus SPL dengan aturan Cramer. & 2 x 2 x 50' tatap muka dan latihan terstruktur; 2 x 2 x 50' tugas mandiri (worksheet). & Mahasiswa menyelesaikan masalah SPL dengan aturan Cramer. & Kriteria: ketepatan dan penguasaan. Bentuk penilaian: ketepatan jawaban tugas. & Mampu menyelesaikan masalah SPL dengan aturan Cramer. & 2 \\
10 & SCPMK0902-01103 & Vektor dan operasi dot product. & Bentuk: kuliah luring. Metode: Self-Directed Learning. Penugasan: Latihan 1 menyelesaikan soal vektor; Worksheet 4 penyelesaian masalah sederhana dengan konsep vektor. & 2 x 2 x 50' tatap muka dan latihan terstruktur; 2 x 2 x 50' tugas mandiri (worksheet). & Mahasiswa memahami konsep vektor, memodelkan permasalahan dalam bentuk vektor, melakukan operasi dot product, dan mengetahui pemanfaatan vektor. & Kriteria: ketepatan dan penguasaan. Bentuk penilaian: ketepatan jawaban tugas. & Memahami konsep vektor; mampu membuat model vektor dari sebuah permasalahan; mampu melakukan operasi dot product; mengetahui pemanfaatan vektor pada bidang informatika. & 1 \\
11 & SCPMK0902-01103 & Operasi vektor cross product. & Bentuk: kuliah luring. Metode: Self-Directed Learning. Penugasan: Latihan 1 dan Worksheet 4 menyelesaikan permasalahan vektor dengan operasi cross product. & 2 x 2 x 50' tatap muka dan latihan terstruktur; 2 x 2 x 50' tugas mandiri (worksheet). & Mahasiswa menyelesaikan permasalahan dengan menggunakan operasi cross product pada vektor. & Kriteria: ketepatan dan penguasaan. Bentuk penilaian: ketepatan jawaban tugas. & Mampu menyelesaikan permasalahan vektor dengan menggunakan cross product. & 1 \\
12 & SCPMK0902-01102, SCPMK0902-01103 & Kuis 2: materi pekan 9--11. & Ujian tertulis. & Sesuai jadwal asesmen. & Mahasiswa mengerjakan soal kuis dengan materi pekan 9--11 secara mandiri. & Ujian. & Ketepatan jawaban dengan kunci jawaban. & 10 \\
13 & SCPMK0902-01103 & Eigenvector dan eigenvalue. & Bentuk: kuliah luring. Metode: Self-Directed Learning. Penugasan: Latihan 1 dan Worksheet 4 menyelesaikan permasalahan dengan eigenvector dan eigenvalue. & 2 x 2 x 50' tatap muka dan latihan terstruktur; 2 x 2 x 50' tugas mandiri (worksheet). & Mahasiswa menyelesaikan permasalahan dengan eigenvector dan eigenvalue serta menggunakannya pada masalah bidang informatika. & Kriteria: ketepatan dan penguasaan. Bentuk penilaian: ketepatan jawaban tugas. & Mampu menghitung nilai eigenvalue dan eigenvector; mampu menyelesaikan masalah informatika dengan eigenvector dan eigenvalue. & 1 \\
14 & SCPMK0902-01103 & Proyeksi ortogonal (ortogonalitas). & Bentuk: kuliah luring. Metode: Self-Directed Learning. Penugasan: Latihan 1 dan Worksheet 4 menyelesaikan permasalahan proyeksi ortogonal. & 2 x 2 x 50' tatap muka dan latihan terstruktur; 2 x 2 x 50' tugas mandiri (worksheet). & Mahasiswa melakukan proses proyeksi ortogonal dan memanfaatkan konsep proyeksi ortogonal dalam penyelesaian masalah informatika. & Kriteria: ketepatan dan penguasaan. Bentuk penilaian: ketepatan jawaban tugas. & Mampu melakukan proyeksi ortogonal; mampu menyelesaikan permasalahan dengan menggunakan proyeksi ortogonal. & 1 \\
15--16 & SCPMK0902-01101, SCPMK0902-01102, SCPMK0902-01103 & Studi kasus pemanfaatan aljabar linier dalam bidang informatika. & Bentuk: kuliah luring. Metode: Case Method. Penugasan: Worksheet 5 menyelesaikan studi kasus bidang informatika. & 2 x 2 x 50' tatap muka dan latihan terstruktur; 2 x 2 x 50' tugas mandiri (worksheet). & Mahasiswa mengetahui manfaat aljabar linier dalam informatika dan memanfaatkan konsep aljabar linier pada permasalahan bidang informatika. & Kriteria: ketepatan dan penguasaan. Bentuk penilaian: ketepatan jawaban tugas. & Mampu menyelesaikan masalah bidang informatika yang berhubungan dengan aljabar linier. & 4 \\
17 & SCPMK0902-01101, SCPMK0902-01102, SCPMK0902-01103 & Ujian Akhir Semester: materi pekan 1--16. & Ujian tertulis. & Sesuai jadwal asesmen. & Mahasiswa mengerjakan soal ujian dengan materi pekan 1--16 secara mandiri. & Ujian. & Ketepatan jawaban dengan kunci jawaban. & 35 \\
\end{longtblr}
